\chapter{Diseño e implementación}
\label{ch:implementacion}

% =======================================================================
% TODO GENERAL: este capítulo sustituye a solucion.tex del anteproyecto.
% Gran parte del texto de solucion.tex se reutiliza convertido a voz
% ejecutada (pasado/presente). Lo NUEVO respecto al anteproyecto es el
% esquema de categorías por firma de energía (decisión 2026-06-11), que
% aún NO está implementado en RTL: al implementarlo, documentar aquí la
% versión final.
% =======================================================================

% TODO: párrafo introductorio — los tres componentes (medición, clasificador,
% análisis), reusar apertura de solucion.tex con la figura de arquitectura general.

% -----------------------------------------------------------------------
\section{Definición de categorías por firma de energía}
\label{sec:categorias_energia}

% TODO: esta sección es NUEVA y es el aporte metodológico central de la
% tesis frente al anteproyecto (que usaba Arithmetic/Logic/Memory/Branch/
% Jump/Float por tipo de operación). Documentar:
%
% 1. Principio: las fronteras entre categorías se justifican A PRIORI desde
%    la microarquitectura del RTL (unidad funcional x latencia), no desde
%    la semántica de la instrucción.
% 2. Latencias reales del CV32E40P (verificadas en el RTL):
%    - ALU: 1 ciclo; mul (MUL_MAC32): 1 ciclo
%    - mulh/mulhu/mulhsu: ~4 ciclos (FSM en cv32e40p_mult.sv)
%    - div/divu/rem/remu: variable hasta ~32 ciclos, dependiente de datos
%      (cv32e40p_alu_div.sv)
%    - branch resuelto en EX; flush solo si es tomado (branch_taken_ex)
% 3. Categorías resultantes:
%    (1) ALU simple: aritmética/lógica/shifts/compares/mul de 1 ciclo
%        + branches NO tomados (equivalen a un compare en EX)
%    (2) Multiplicación larga: mulh/mulhu/mulhsu
%    (3) División: div/divu/rem/remu
%    (4) Memoria: loads/stores incluidos flw/fsw (activan data_req)
%    (5) Control: jal/jalr + branches TOMADOS (mhpmevent_jump ||
%        mhpmevent_branch_taken; el flush/refetch es un mecanismo único)
%    (6) Flotante: cómputo FPU
% 4. Instrucciones NO contadas (filtro): csrr/csrw, fence/fence.i, ecall,
%    ebreak, mret, wfi --> sum(n_i) <= minstret. Justificar: el decoder
%    deja alu_en=1/ALU_SLTU por defecto y sin filtro contaminarían la
%    categoría ALU simple.
% 5. Si la caracterización muestra que mulh es despreciable en código real,
%    se pliega a ALU simple y quedan 5 categorías (documentar la decisión).

% -----------------------------------------------------------------------
\section{Módulo clasificador de instrucciones}
\label{sec:clasificador_impl}

% TODO: reusar de solucion.tex §3.2 (concepto, señales del pipeline) en voz
% ejecutada, actualizado al esquema nuevo:
% - entradas: mhpmevent_minstret (retire), mhpmevent_load/store,
%   mhpmevent_jump, mhpmevent_branch_taken (NUEVA), mult_en + selector
%   mulh, señal de división, apu_en, csr_access_ex, system_insn (NUEVA)
% - tabla categoría -> condición de detección (rehacer tab:senales_clas)
% - contadores de 64 bits, lógica de incremento

\subsection{Interfaz CSR}
\label{sec:csr_impl}

% TODO: 12 CSR en pares LO/HI, rango 0xBC0--0xBCB (custom RW de modo
% máquina). Tabla de mapeo. Lectura (HI<<32)|LO y reseteo por escritura.
% Si el esquema nuevo cambia el mapeo, actualizar también COMANDOS_JTAG.

\subsection{Integración en el núcleo}
\label{sec:integracion_impl}

% TODO: reusar solucion.tex §3.2.3: instancia en cv32e40p_core.sv,
% multiplexor de lectura csr_rdata, decoder aceptando los CSR custom,
% manifiestos (Bender.yml, flist, src_files.yml).

% -----------------------------------------------------------------------
\section{Plataforma FPGA}
\label{sec:plataforma_impl}

% TODO: nueva sección (en el anteproyecto estaba disperso). Documentar lo
% ya hecho: target pulpissimo-nexys, pad_esp32_sync (io_08 -> H4/PMOD JD1),
% fix de dirección de pad_jtag_tdo, JTAG por PMOD JA con FT232H,
% constraints de temporización JTAG, síntesis no incremental
% (fix Synth 8-6885), frecuencia de operación verificada (10 MHz).

% -----------------------------------------------------------------------
\section{Circuito de medición de potencia}
\label{sec:medicion_impl}

% TODO: reusar solucion.tex §3.1 completo en voz ejecutada (INA240A1 +
% shunt 100 mOhm high-side, ADS1115, ESP32, resolución 156 uW, ventana por
% GPIO de sincronización). Añadir fotos/esquemático del hardware REAL
% construido y cualquier desviación respecto a lo propuesto.

% -----------------------------------------------------------------------
\section{Firmware de caracterización}
\label{sec:firmware_impl}

% TODO: bucles dominados por categoría (~97% dominancia, 64+2 instrucciones
% por iteración), delimitación por GPIO8, terminación con ebreak,
% frequency_probe. Actualizar las rutinas al esquema de categorías nuevo
% (p. ej. rutinas separadas para división y mulh; branches tomados vs no
% tomados).

% -----------------------------------------------------------------------
\section{Software de análisis y operación}
\label{sec:software_impl}

% TODO: reusar solucion.tex §3.3 en voz ejecutada: calibración (M1/M2),
% validación y modo de operación sin instrumentación. Documentar los
% scripts reales y el formato de coefficients.json.

%%% Local Variables:
%%% mode: latex
%%% TeX-master: "main"
%%% End:
